% Part: first-order-logic
% Chapter: syntax-and-semantics
% Section: free-vars-sentences

\documentclass[../../../include/open-logic-section]{subfiles}

\begin{document}

\olfileid[fr]{fol}{syn}{fvs}

\olsection{\printtoken{P}{variable} libres et \printtoken{p}{sentence}}

\begin{defn}[Occurrences libres d'!!a{variable}]
\ollabel{defn:free-occ}
Les occurrences \emph{libres} d'!!a{variable} dans !!a{formula}
sont définies inductivement comme suit :
\begin{enumerate}
\item \indcase*{!A}{$!A$ est atomique}{toutes les occurrences de
  !!{variable} dans $\indfrm$ sont libres.}

\tagitem{prvNot}{\indcase{!A}{\lnot !B}{les occurrences libres de
  !!{variable} dans $\indfrm$ sont exactement celles de~$!B$.}}{}

\item \indcase{!A}{(!B \ast !C)}{les occurrences libres de
  !!{variable} dans $\indfrm$ sont celles de~$!B$ réunies à
  celles de~$!C$.}

\tagitem{prvAll}{\indcase{!A}{\lforall[x][!B]}{les occurrences
  libres de !!{variable} dans $\indfrm$ sont toutes celles de~$!B$,
  sauf les occurrences de~$x$.}}{}

\tagitem{prvEx}{\indcase{!A}{\lexists[x][!B]}{les occurrences
  libres de !!{variable} dans $\indfrm$ sont toutes celles de~$!B$,
  sauf les occurrences de~$x$.}}{}
\end{enumerate}
\end{defn}

\begin{defn}[Occurrences liées de variables]
Une occurrence d'!!a{variable} dans une !!{formula}~$!A$ est
\emph{liée} si elle n'est pas libre.
\end{defn}

\begin{prob}
Donner une définition inductive des occurrences liées de variables,
sur le modèle de la \olref[fol][syn][fvs]{defn:free-occ}.
\end{prob}

\begin{defn}[Portée]
\iftag{prvAll}{Si $\lforall[x][!B]$ apparaît comme occurrence d'une
  sous-formule dans une !!{formula}~$!A$, l'occurrence correspondante
  de~$!B$ dans~$!A$ est la \emph{portée} de l'occurrence correspondante
  de~$\lforall[x]$. \iftag{prvEx}{Il en va de même pour
  $\lexists[x]$.}{}}{Si $\lexists[x][!B]$ apparaît comme occurrence
  d'une sous-formule dans une !!{formula}~$!A$, l'occurrence
  correspondante de~$!B$ dans~$!A$ est la \emph{portée} de
  l'occurrence correspondante de~$\lexists[x]$.}

Si $!B$ est la portée d'une occurrence du quantificateur
\iftag{prvAll}{$\lforall[x]$\iftag{prvEx}{ ou
    $\lexists[x]$}{}}{$\lexists[x]$} dans~$!A$, les occurrences
libres de~$x$ dans~$!B$ sont liées dans
\iftag{prvAll}{$\lforall[x][!B]$\iftag{prvEx}{ et
$\lexists[x][!B]$}{}}{$\lexists[x][!B]$}. Nous disons que ces
occurrences sont \emph{liées par} l'occurrence du quantificateur en
question.
\end{defn}

\begin{ex}
Considérons la !!{formula} suivante :
\[
\lexists[\Obj v_0][\underbrace{\Atom{\Obj A^2_0}{\Obj v_0,\Obj v_1}}_{!B}]
\]
$!B$ représente la portée de~$\lexists[\Obj v_0]$. Ce
quantificateur lie l'occurrence de~$\Obj v_0$ dans~$!B$, mais non
celle de~$\Obj v_1$. Ainsi, $\Obj v_1$ est ici une variable libre.

Voyons maintenant une !!{formula}~$!A$ plus complexe :
\[
\lforall[\Obj v_0][\underbrace{(\Atom{\Obj A^1_0}{\Obj v_0} \lif
    \Atom{\Obj A^2_0}{\Obj v_0, \Obj v_1})}_{!B}] \lif \lexists[\Obj
  v_1][\underbrace{(\Atom{\Obj A^2_1}{\Obj v_0, \Obj v_1} \lor \lforall[\Obj v_0][\overbrace{\lnot \Atom{\Obj A^1_1}{\Obj v_0}}^{!D}])}_{!C}]
\]
$!B$ est la portée du premier $\lforall[\Obj v_0]$ ; $!C$ est
la portée de $\lexists[\Obj v_1]$ et $!D$ celle du second
$\lforall[\Obj v_0]$. Le premier $\lforall[\Obj v_0]$ lie les
occurrences de~$\Obj v_0$ dans~$!B$ ; $\lexists[\Obj v_1]$ lie
l'occurrence de~$\Obj v_1$ dans~$!C$ ; le second
$\lforall[\Obj v_0]$ lie l'occurrence de~$\Obj v_0$ dans~$!D$.
La première occurrence de~$\Obj v_1$ et la quatrième occurrence
de~$\Obj v_0$ sont libres dans~$!A$. La dernière occurrence
de~$\Obj v_0$ est libre dans~$!D$, mais liée dans~$!C$ et~$!A$.
\end{ex}

\begin{defn}[Énoncé]
Une !!{formula}~$!A$ est un \emph{!!{sentence}} si et seulement si
elle ne contient aucune occurrence libre de !!{variable}.
\end{defn}

% Ajouter des exemples.

\end{document}
